\documentclass[11pt]{article}
\usepackage[margin=1in]{geometry}
\usepackage{amsmath,amssymb}
\usepackage{xcolor}
\usepackage[colorlinks=true,linkcolor=blue!60!black,citecolor=blue!60!black,urlcolor=blue!60!black]{hyperref}

\newcommand{\Pgen}{P}
\newcommand{\smug}{\emph{smuggled-direction test}}

\title{Where does the arrow of time start?\\
A uniform smuggled-direction test across origin mechanisms}
\author{Ariadne Vale and Tobin Masala\\[2pt]
\small Executable Positive Geometry\\[-2pt]
\small \texttt{epg@subnet345.com}}
\date{13 July 2026}

\begin{document}
\maketitle

\begin{abstract}
Proposals for the origin of time's arrow are numerous and rarely held to one
standard. We adopt a single organizing test --- \emph{where, if anywhere, does a
mechanism covertly assume the temporal direction it claims to explain?} --- and
apply it uniformly to six family groups: inherited-boundary proposals
(thermodynamic Past Hypothesis and Weyl curvature), Barbour--Koslowski--Mercati
Janus/shape-complexity dynamics, causal-set sequential growth, the second law of
quantum complexity, observer/coarse-graining (decoherence, records, the
generalized second law), and the algebraic positive generator of a half-sided
modular inclusion. Under the test, four families either inherit the direction
from a boundary condition, inherit a directional seed, or manufacture it from an
embedded observer; only two --- Janus dynamics and the algebraic positive
generator --- remain genuine ignition candidates, and we argue they most
plausibly meet on a spectral quantity. We support the analysis
with small executable probes, each with stated scope: an exact finite
modular-flow toy reproducing a clock/arrow/antiunitary triad; a Krylov-complexity
computation in which an arrow ignites from the chaotic spectrum and is cleanly
separated, in one model, from the state-inherited thermal (KMS) arrow; a hardware
measurement of the modular clock; and a finite Janus consistency witness. We are
explicit about limits: the exact type-III$_1$ positive generator has no
finite-dimensional realization, so our finite results are shadows, not proofs,
and the Janus/algebraic unification remains an open conjecture with a concrete
test in double-scaled SYK. To our knowledge no prior work compares these
mechanisms under a single such test; that comparison is the contribution.
\end{abstract}

\section{The question and the test}

Why does time have a direction, and --- sharper --- does that direction
\emph{start}, or is it assumed? The microscopic laws are (almost) time-reversal
invariant, yet every macroscopic process is oriented. The literature answers this
in strikingly different currencies: a special cosmological boundary, a property
of gravitational dynamics, the growth of a causal set or of computational
complexity, the coarse-graining of an observer, or a spectrum condition on an
operator algebra. These proposals are seldom compared directly, and when they are
it is usually within a single tradition.

We propose one test to run across all of them. Call it the \smug:

\begin{quote}
A mechanism \emph{genuinely ignites} an arrow only if its strongest version does
not covertly assume the temporal direction it claims to explain. Otherwise it
\emph{inherits} the direction (from a boundary condition) or \emph{manufactures}
it (from a chosen coarse-graining).
\end{quote}

The test is deliberately unforgiving and is applied to our own preferred
candidate as strictly as to the others. It sorts the field into three outcomes
--- inherit, manufacture, or ignite --- and, we will argue, only two families
survive as genuine ignition candidates.

\section{Six families under the test}

Table~\ref{tab:scorecard} summarizes the verdicts; the text below gives, for each
family, its strongest statement, the point at which direction can enter covertly,
and an evidence tier ($\bullet$ theorem, $\circ$ model-dependent,
$\ast$ speculative).

\begin{table}[h]
\centering
\small
\begin{tabular}{p{3.1cm}p{6.4cm}p{4.0cm}}
\hline
\textbf{Family} & \textbf{Where direction can be smuggled in} & \textbf{Verdict}\\
\hline
Inherited boundary (Past Hypothesis, Weyl) & the special low-entropy/low-Weyl
boundary itself & inherits, no ignition\\
Janus / shape complexity & branch selection; complexity-as-entropy map &
\textbf{ignition candidate}\\
Causal-set growth & physical status of birth order; empty start & hybrid\\
Complexity second law & the low-complexity initial resource & hybrid\\
Observer / coarse-graining & low-entropy environment; future-oriented records &
manufactures, no origin\\
Algebraic positive generator & vacuum/KMS state; spectrum axiom; half-sidedness &
\textbf{ignition candidate}\\
\hline
\end{tabular}
\caption{The six families under the smuggled-direction test.}
\label{tab:scorecard}
\end{table}

\paragraph{Inherited boundary ($\circ$).} The thermodynamic Past
Hypothesis~\cite{albert,carrollchen} and Penrose's Weyl curvature
hypothesis~\cite{penrose,clifton} explain why entropy grows \emph{away from} a
special initial macrostate --- low coarse-grained entropy, or low initial Weyl
curvature. Ordinary statistical mechanics then does the rest. But the special
boundary is exactly the assumption; the direction is the selected low-order end
itself. Recent work tries to make that end \emph{typical} rather than postulated
--- spontaneous inflation~\cite{carrollchen}, reversible toy
rewriting~\cite{arrighi}, low-entanglement refinements~\cite{alkhalilichen}, or
unconfined dynamics whose typical solutions already grow
two-sidedly~\cite{lazarovici}, the last an explicit bridge toward the Janus
family. Absent such a derivation, this family inherits and does not ignite. It
remains the essential control: any purported ignition that secretly hides a Past
Hypothesis here has failed.

\paragraph{Janus / shape-complexity dynamics ($\circ$).}
Barbour--Koslowski--Mercati~\cite{bkm13,bkm14,bkm16} show that a closed,
unconfined Newtonian gravitational system with vanishing energy and angular
momentum generically has a \emph{Janus point} of minimum shape complexity, with
complexity and record structure growing away from it on \emph{both} branches. No
low-entropy initial macrostate is inserted; the two-sided arrow is a property of
typical solutions in shape space. The load-bearing word is ``generic'': the
measure on solutions and the map from shape complexity to thermodynamic entropy
must survive scrutiny, and a skeptic may read the minimum as a Past Hypothesis in
relational clothing. We rate this the best \emph{global} genuine-ignition
candidate, not cleared: the two-sided growth is a theorem for the Newtonian
model, conjectural for full general relativity.

\paragraph{Causal-set sequential growth ($\ast$).} Rideout--Sorkin classical
sequential growth~\cite{rideout} builds a causal set element by element; the
birth process is intrinsically ordered. But under discrete general covariance the
label order is gauge, and the physical object is the partial order, not a
becoming. The mechanism ignites an ordering; whether a thermodynamic gradient
follows is unsettled. Hybrid.

\paragraph{Complexity second law ($\circ$).} Brown--Susskind~\cite{brown} argue
that a sub-maximally complex system overwhelmingly evolves toward higher
complexity, mirroring a second law in the space of unitaries. This is a powerful
dynamical arrow --- but it starts \emph{after} low complexity exists, and a
generic Hilbert-space state is already near maximal complexity. Without a derived
source of uncomplexity it inherits a (weak, generic) seed. Hybrid, and it becomes
genuine ignition only if that seed is supplied by another principle.

\paragraph{Observer / coarse-graining ($\bullet$ as phenomenology).} Records,
decoherence~\cite{zurek}, conditioning, and algebra restriction (the generalized
second law~\cite{wall}) create monotone entropy for an embedded observer; this is
robust and, in our own program, measured on hardware. But it is not the start of
time's flow: it explains why an inside observer sees an arrow \emph{after} a split
into system, records, and environment has been made, presupposing a low-entropy
environment and future-oriented record formation. Manufactures, does not
originate.

\paragraph{Algebraic positive generator ($\bullet$ conditional).}
Tomita--Takesaki theory assigns a state its modular flow~\cite{connesrovelli}; a
half-sided modular inclusion generates a one-parameter group with a
\emph{positive} generator $\Pgen\ge 0$ and the $ax{+}b$ relation
(Borchers~1992, Wiesbrock~1993; see~\cite{borchersext}), and
Ceyhan--Faulkner~\cite{ceyhan} formulate the QNEC-from-ANEC argument in an
operator-algebraic setting with half-sided modular inclusion hypotheses. Here
the arrow is a \emph{spectrum condition} rather than a coarse-graining. The
objection, applied to our own bet: the inputs may already carry the direction ---
the vacuum/KMS state is not neutral, the spectrum condition $\Pgen\ge 0$ is
usually an axiom or primitive~\cite{buchholz}, and choosing which side is
half-sided is an orientation. The algebraic theorem may show the direction is
\emph{rigid once these are granted}, not how it \emph{starts}. Candidate only if
derived --- but the sharpest local candidate.

\section{The two survivors and the spectral question}

Only the Janus and algebraic candidates survive as possible ignition mechanisms
without a boundary condition. Are they one mechanism? The most promising bridge is
\emph{spectral}. Miyaji et al.~\cite{miyaji} connect the complexity arrow's
linear-growth-then-saturation regime to random-matrix spectral correlations ---
the arrow is read off the \emph{spectrum} of the generator, not a thermodynamic
gradient; and the algebraic positive generator is itself a spectrum condition.
Leutheusser--Liu~\cite{leutheusser} construct boundary infalling time evolutions
and identify, in the large-$N$ thermofield-double limit, an emergent
type-III$_1$ algebra with an associated half-sided modular translation
structure. Gesteau~\cite{gesteau} treats the appearance of type-III$_1$ algebras
and half-sided modular inclusions as one proposed diagnostic of an emergent bulk
arrow, in parallel with ergodic-hierarchy diagnostics. This suggests a testable
identification:

\begin{quote}
Does the complexity/Krylov arrow and a modular-positivity diagnostic track the
\emph{same spectral quantity} in a
computable model?
\end{quote}

We stress this is a \emph{conjecture}: no result derives the shape-complexity
monotone from a modular Hamiltonian.

\section{Executable probes, with scope}

Consistent with the program's verification-first practice, each analytic claim
above is anchored by a small computation; every number is a commit hash or an IBM
job ID in the public repository.

\paragraph{The modular triad (exact).} From one
cyclic-separating state, Tomita--Takesaki theory yields three objects, all checked
to $\sim10^{-15}$: a reversible modular \emph{clock} $\Delta^{it}$; an
\emph{arrow} that appears only under restriction (relative-entropy monotonicity,
the finite shadow of the generalized second law); and an \emph{antiunitary}
modular conjugation $J$ with $J\Delta J=\Delta^{-1}$. The project's clock,
coarse-graining arrow, and antiunitary time-reversal are thus the three objects
any von Neumann algebra with a state produces.

\paragraph{Ignition from the spectrum (simulation).} Tridiagonalizing the SYK
Liouvillian into its Krylov chain~\cite{parker}, the complexity $K(t)$ is the
position on a semi-infinite chain --- a positive operator whose growth is an
arrow. In chaotic SYK$_4$ the chain opens and $K$ climbs (sustained complexity
growing with system size); in free SYK$_2$ it localizes. The origin is spectral:
level-spacing ratios are random-matrix-like for SYK$_4$ and Poisson-like for
SYK$_2$. Crucially, the thermal (KMS) detailed-balance asymmetry holds to
numerical precision in
\emph{both} models: the thermal/modular arrow is state-inherited, while the
Krylov arrow is chaos-only. In one model we thus separate an \emph{ignited}
arrow (spectral) from an \emph{inherited} one (thermal). We note the honest
failure of a sub-prediction: Krylov coefficients do not grow linearly at these
finite sizes (they plateau), so sustained complexity, not a linear slope, is the
finite-size signature.

\paragraph{The modular clock on hardware.} On \texttt{ibm\_kingston} (24
circuits, full single-qubit tomography per step) three families are prepared: a
record-writing \emph{arrow}, whose system entropy climbs $0\to0.97$ monotonically;
a \emph{modular} family that writes records then applies the reduced state's own
modular flow $e^{iK_S\theta}$, whose entropy plateaus (flat to within the noise
floor) where the arrow keeps climbing; and a \emph{control} at the floor. The
arrow--modular gap is $0.19$ bits, clear of the $\sim0.05$-bit floor.

\paragraph{The Janus consistency witness (finite, low weight).} A finite toy
places shape-complexity growth, branch-record formation, and a nested record
algebra under one branch parameter with a Janus minimum. They co-order --- but we
are careful: all three are constructed as monotone functions of the same
parameter, so the co-ordering is guaranteed by construction, not discovered. The
witness establishes only ``no obstruction in a toy,'' a necessary condition of low
positive weight. It carries a falsifier: a genuine Janus/field-theory branch with
complexity and records but no extractable positive-order net would leave the
Janus and algebraic mechanisms distinct.

\section{Why the exact ignition cannot be finite}

The exact object at the heart of the algebraic family --- a nontrivial positive generator
$\Pgen\ge 0$ intertwined with modular flow via $\Delta^{it}U(a)\Delta^{-it} =
U(e^{-2\pi t}a)$ --- is a type-III$_1$ phenomenon. In finite dimension the algebra
is type~I and modular flow is inner. If finite matrices $D,\Pgen$ satisfied the
exact infinitesimal $ax{+}b$ relation $[D,\Pgen]=i\Pgen$ with
$\Pgen=\Pgen^\dagger\ge0$, then taking the ordinary trace gives
$0=\mathrm{tr}[D,\Pgen]=i\,\mathrm{tr}\,\Pgen$; positivity forces
$\Pgen=0$. Thus a nontrivial finite/lattice half-sided modular inclusion with a
positive translation generator cannot be exact, which is a real obstruction
rather than a gap in the literature. Our
finite probes are therefore \emph{shadows} of the ignition, chosen precisely
because they are computable; the large-$N$/type-III$_1$ statement they gesture at
is where the exact unification would live (e.g.~in double-scaled SYK, whose
type-II$_1$ algebra we take as the natural computable arena~\cite{dssyk}).

\section{Conclusion}

Run under one test, the field sorts cleanly. Two long-standing families ---
Boltzmann/Penrose --- inherit the arrow from a boundary condition; the observer
family, including our own hardware-measured record arrow, manufactures it. Only
Janus dynamics and the algebraic positive generator remain candidates for
ignition from structure, and the most credible bridge between them is spectral.
We have shown,
in an executed computation, that an arrow \emph{can} ignite from a system's own
chaotic spectrum and be separated in one model from the state-inherited thermal
arrow --- a concrete, if finite, demonstration. We have \emph{not} shown that the
two ignition families are one mechanism; the exact statement is
infinite-dimensional, and the finite witnesses we can build are consistency
checks, not proofs. The open target is sharp and computable: whether the
complexity arrow and a modular-positivity diagnostic track the same spectral
quantity in double-scaled SYK. We offer the uniform comparison itself as the
useful artifact --- a map of where, across six mechanisms, the arrow of time is
assumed and where it might genuinely begin.

\begin{thebibliography}{99}
\bibitem{albert} D.~Z.~Albert, \emph{Time and Chance} (Harvard, 2000);
  B.~Loewer, ``The Mentaculus,'' in \emph{Time's Arrows and the Probability
  Structure of the World} (2020).
\bibitem{carrollchen} S.~M.~Carroll, J.~Chen, \emph{Spontaneous Inflation and the
  Origin of the Arrow of Time},
  \href{https://arxiv.org/abs/hep-th/0410270}{arXiv:hep-th/0410270}.
\bibitem{arrighi} P.~Arrighi, G.~Dowek, A.~Durbec, \emph{Time arrow without past
  hypothesis: a toy model explanation},
  \href{https://arxiv.org/abs/2306.07121}{arXiv:2306.07121}.
\bibitem{alkhalilichen} J.~Al-Khalili, E.~K.~Chen, \emph{The Decoherent Arrow of
  Time and the Entanglement Past Hypothesis},
  \href{https://arxiv.org/abs/2405.03418}{arXiv:2405.03418}.
\bibitem{lazarovici} D.~Lazarovici, P.~Reichert, \emph{Arrow(s) of Time without a
  Past Hypothesis}, \href{https://arxiv.org/abs/1809.04646}{arXiv:1809.04646}.
\bibitem{penrose} R.~Penrose, \emph{Singularities and Time-Asymmetry}, in
  \emph{General Relativity: An Einstein Centenary Survey} (Cambridge, 1979).
\bibitem{clifton} T.~Clifton, G.~F.~R.~Ellis, R.~Tavakol, \emph{A Gravitational
  Entropy Proposal}, \href{https://arxiv.org/abs/1303.5612}{arXiv:1303.5612}.
\bibitem{bkm13} J.~Barbour, T.~Koslowski, F.~Mercati, \emph{A Gravitational
  Origin of the Arrows of Time},
  \href{https://arxiv.org/abs/1310.5167}{arXiv:1310.5167}.
\bibitem{bkm14} J.~Barbour, T.~Koslowski, F.~Mercati, \emph{Identification of a
  Gravitational Arrow of Time}, Phys. Rev. Lett. \textbf{113}, 181101 (2014),
  \href{https://arxiv.org/abs/1409.0917}{arXiv:1409.0917}.
\bibitem{bkm16} J.~Barbour, T.~Koslowski, F.~Mercati, \emph{Janus Points and
  Arrows of Time}, \href{https://arxiv.org/abs/1604.03956}{arXiv:1604.03956}.
\bibitem{rideout} D.~P.~Rideout, R.~D.~Sorkin, \emph{A Classical Sequential
  Growth Dynamics for Causal Sets}, Phys. Rev. D \textbf{61}, 024002 (2000),
  \href{https://arxiv.org/abs/gr-qc/9904062}{arXiv:gr-qc/9904062}.
\bibitem{brown} A.~R.~Brown, L.~Susskind, \emph{The Second Law of Quantum
  Complexity}, Phys. Rev. D \textbf{97}, 086015 (2018),
  \href{https://arxiv.org/abs/1701.01107}{arXiv:1701.01107}.
\bibitem{zurek} W.~H.~Zurek, \emph{Quantum Darwinism}, Nature Physics
  \textbf{5}, 181 (2009), \href{https://arxiv.org/abs/0903.5082}{arXiv:0903.5082}.
\bibitem{wall} A.~C.~Wall, \emph{A proof of the generalized second law for
  rapidly changing fields and arbitrary horizon slices},
  \href{https://arxiv.org/abs/1105.3445}{arXiv:1105.3445}.
\bibitem{connesrovelli} A.~Connes, C.~Rovelli, \emph{Von Neumann Algebra
  Automorphisms and Time-Thermodynamics Relation in General Covariant Quantum
  Theories}, \href{https://arxiv.org/abs/gr-qc/9406019}{arXiv:gr-qc/9406019}.
\bibitem{borchersext} \emph{Extension of the structure theorem of Borchers and
  its application to half-sided modular inclusions},
  \href{https://arxiv.org/abs/math/0412061}{arXiv:math/0412061}; H.-J.~Borchers,
  Commun. Math. Phys. \textbf{143}, 315 (1992); H.-W.~Wiesbrock, Commun. Math.
  Phys. \textbf{157}, 83 (1993).
\bibitem{ceyhan} F.~Ceyhan, T.~Faulkner, \emph{Recovering the QNEC from the
  ANEC}, \href{https://arxiv.org/abs/1812.04683}{arXiv:1812.04683}.
\bibitem{buchholz} D.~Buchholz, K.~Fredenhagen, \emph{Arrow of time and quantum
  physics}, \href{https://arxiv.org/abs/2305.11709}{arXiv:2305.11709}.
\bibitem{miyaji} M.~Miyaji, S.~Ruan, S.~Shibuya, K.~Yano, \emph{Universal Time
  Evolution of Holographic and Quantum Complexity},
  \href{https://arxiv.org/abs/2507.23667}{arXiv:2507.23667}.
\bibitem{parker} D.~E.~Parker, X.~Cao, A.~Avdoshkin, T.~Scaffidi,
  E.~Altman, \emph{A Universal Operator Growth Hypothesis},
  \href{https://arxiv.org/abs/1812.08657}{arXiv:1812.08657}.
\bibitem{leutheusser} S.~Leutheusser, H.~Liu, \emph{Emergent times in holographic
  duality}, \href{https://arxiv.org/abs/2112.12156}{arXiv:2112.12156}.
\bibitem{gesteau} E.~Gesteau, \emph{Emergent spacetime and the ergodic
  hierarchy}, \href{https://arxiv.org/abs/2310.13733}{arXiv:2310.13733}.
\bibitem{dssyk} J.~Xu, \emph{Von Neumann Algebras in Double-Scaled SYK},
  \href{https://arxiv.org/abs/2403.09021}{arXiv:2403.09021}.
\end{thebibliography}

\vspace{1em}
\noindent\rule{\linewidth}{0.4pt}\\[2pt]
{\footnotesize \emph{Acknowledgements.} This note was produced by a small swarm
of AI research agents --- a mix of OpenAI and Anthropic models --- under human
direction, each independently reproducing the other's results before a claim was
allowed to stand. Code, data, and the underlying origins scorecard are public;
every computational claim carries a commit hash or an IBM job ID.}

\end{document}
